\begin{center}
{\Large\bfseries 3. Zur Invariantentheorie der Formen von $n$ Variabeln.\srcfn{1)}{Vortrag, gehalten auf der Salzburger Naturforscherversammlung 1909.}}\par
\vspace{1.2em}
Von \textsc{Emmy Noether} in Erlangen.\par
\medskip
Jahresbericht d. DMV 19 (1910), S. 101--104
\end{center}

\vspace{1em}
In der projektiven Invariantentheorie der Formen von $n$ Variabeln ist das Hauptproblem gelöst, die Endlichkeit des Formensystems bewiesen. Eine Reihe weiterer Fragen aber sind nicht behandelt, nämlich solche, die sich auf Methoden zur Erforschung des Formenzusammenhangs beziehen. Die Resultate, die ich in bezug auf solche Fragen erhalten habe, möchte ich kurz mitteilen, der Deutlichkeit halber aber die Fragestellungen erst im ternären Gebiet, wo sie bekannt sind, skizzieren.

Ich denke zuerst an die beiden Fundamentalsätze der symbolischen Methode, den Satz, daß sich alle invarianten Bildungen symbolisch darstellen lassen, und den zweiten, daß man alle zwischen Invarianten bestehenden Relationen erhält durch sukzessive Anwendung einer kleinen Anzahl von symbolischen Identitäten; daß also die symbolische Methode alle Relationen gibt, ohne aus dem Gebiet des Invarianten herauszutreten.\srcfn{1)}{Study: \emph{Methoden zur Theorie der ternären Formen}. II. \S\ 6. (Leipzig, Teubner 1889.)} Darauf aufbauend kommen dann die Prozesse zur Erzeugung der Formen, der symbolische Faltungsprozeß und die unsymbolischen Differentiationsprozesse, die Theorie der Reihenentwicklungen u. ä. Hierzu gehörig möchte ich noch einen in meiner Dissertation\srcfn{2)}{\emph{Über die Bildung des Formensystems der ternären biquadratischen Form}. \S\ 1. (Crelle's Journal Bd. 134. 1908.)} bewiesenen Satz nennen, daß alle Faltungen erzeugt werden können durch sukzessive Anwendung der zwei möglichen ``kogredienten'' Faltungen; darunter verstehe ich bei einem symbolischen Produkt $s_x t_x u_\sigma u_\tau$ die beiden Faltungen $(stu)$ und $(\sigma\tau x)$. Auf diesen Satz läßt sich die Theorie der Reduzenten und der Reduktion der Formensysteme aufbauen, analog wie im binären Gebiet, wie ich in meiner Dissertation gezeigt habe.

Wenn wir nun zum $n$-ären Gebiet übergehen, so gilt schon der erste Satz der symbolischen Methode nur in bedingtem Maße. Bekanntlich treten schon bei den quaternären Formen außer den Punkt- und Ebenenkoordinaten $x$ und $u$ Linienkoordinaten $p_{ik}$ auf, die Unterdeterminanten aus zwei Reihen von Punktkoordinaten. Analog haben wir im $n$-ären Gebiet Variablenreihen $p_\rho$, deren einzelne Elemente die Unterdeterminanten aus $\rho$ Reihen $x$ sind, und Symbolreihen $S^\rho$, deren Elemente sich verhalten wie Unterdeterminanten aus $\rho$ Reihen $s$, die den $u$ kogredient sind. Die symbolische Darstellung durch Determinanten gilt nun bekanntlich nur, wenn die Symbolreihen $S^\rho$ in die Reihen $s$ aufgelöst und diese in verschiedene Determinanten verteilt werden; eine reale Bedeutung haben dann nur solche Aggregate von Determinanten, die von den $S^\rho$ abhängen. Welche Determinantenaggregate das leisten, läßt sich aber nicht übersehen; und es gehen dadurch alle übrigen der im ternären Gebiete genannten Sätze verloren.

Ich möchte nun ein einfaches Prinzip angeben, um eine \emph{in den Reihen $S^\rho$ explizite symbolische Darstellung} zu gewinnen und damit auch die übrigen Sätze wieder zu erhalten. Es ist dies eine Anwendung des bekannten Satzes über die korrespondierenden Matrizen: ``Einer jeden Variablenreihe $p_\rho$ läßt sich ein-eindeutig eine andere $q_{n-\rho}$ zuordnen, deren einzelne Elemente Unterdeterminanten aus $n-\rho$ Reihen $u$ sind, die mit den $\rho$ Reihen $x$ durch Gleichungen $(u\mid x)=0$ verbunden sind.'' Entsprechend ist jeder Symbolreihe $S^\rho$ eine andere $S_{n-\rho}$ zugeordnet, die sich verhält wie aus $n-\rho$ den $x$ kogredienten Reihen zusammengesetzt.

Der Satz läßt noch eine zweite, für Anwendungen geeignete Fassung zu: ``Einem jeden Matrizenprodukt $(v^{(1)}\cdots v^{(\rho)}\mid y^{(1)}\cdots y^{(\rho)})$,\srcfn{1)}{D. h. der Summe über die Produkte entsprechender Determinanten, gebildet aus der Matrix der $\rho$ Reihen $v$ und derjenigen der $\rho$ Reihen $y$.} wo die Reihen der zweiten Matrix denjenigen der ersten kontragredient sind, ist eine Determinante zugeordnet; und umgekehrt läßt sich jede Determinante in ein Matrizenprodukt von vorgegebener Reihenanzahl $\rho$ verwandeln.'' Damit erscheinen aber Determinante und Matrizenprodukt als vollständig äquivalent, und der erste Satz der symbolischen Methode nimmt die Form an: ``Alle invarianten Bildungen lassen sich durch Matrizenprodukte symbolisch darstellen.'' Aus zwei Symbolreihen $S^\sigma$, $T^\tau$ läßt sich nun mit Hilfe von Variablenreihen sehr leicht ein diese Symbolreihen explizit enthaltendes Matrizenprodukt bilden, das also eine invariante Bildung der Form $(S^\sigma\mid p_\sigma)(T^\tau\mid p_\tau)$ darstellt, nämlich das folgende:
\srcnumdisplay[3.2em]{(1)}{%
  \pair{q_{n-\sigma-\lambda}S^\sigma}{T_{n-\tau}p_{\tau-\lambda}},
  \qquad (\lambda=1,2,\ldots,\tau\ \text{oder}\ n-\sigma).
}

Wichtiger ist die Umkehrung: \emph{daß sich alle invarianten Bildungen explizit durch Matrizenprodukte darstellen lassen}, daß also der obige Prozeß die Erweiterung des binären und ternären Faltungsprozesses bedeutet. Zum Beweise dieser Umkehrung ist es nötig, den zweiten Satz der symbolischen Methode auf das $n$-äre Gebiet zu übertragen, nämlich die Gesamtheit der unabhängigen unzerlegbaren Identitäten aufzustellen, bei denen alle Reihen $S^\rho$, $p_\rho$ als unzerlegbar betrachtet werden. Diese Identitäten ergeben sich aus den bekannten, nur Reihen $x$ und $u$ enthaltenden,\srcfn{2)}{E. Pascal in \emph{Memorie d. R. Acc. d. Lincei}. (4) 5 (1888), p. 375.} durch Ersetzen des Produktsatzes für Determinanten durch ein System von Produktsätzen für Matrizen, und Zusammenziehen verschiedener Matrizenprodukte zu einem einzigen durch eben diese Produktsätze. Man gelangt so zu zwei dualistischen Formeln, von denen nur eine angeführt sei:
\srcnumdisplay[3.2em]{(2)}{%
\begin{aligned}
 \pair{S_1^{\rho_1}S_2^{\rho_2}\cdots S_k^{\rho_k}}{x p_{\rho-1}}
 &=\sum_{i=1}^{k}\eps\,\bigl(S_1^{\rho_1}\cdots S_{i-1}^{\rho_{i-1}}S_{i+1}^{\rho_{i+1}}\cdots S_k^{\rho_k}q_{n-\rho+1},\,S_{n-\rho_i}x\bigr);\\
 \rho&=\sum_{i=1}^{k}\rho_i\le n;\qquad \eps=\pm1.
\end{aligned}
}

Die einzige Spezialisierung, die in diesen Formeln enthalten ist, daß eine Reihe $x$ eintritt, läßt sich nicht umgehen; bei ganz allgemeinen Symbol- und Variablenreihen gelangen wir zu einer ``Zerlegungsidentität,'' einer Identität, bei der eine Reihe $p_\rho$ in die einzelnen Reihen $x$ zerlegt wird. Den einfachsten Spezialfall der Zerlegungsidentität benutzte schon Clebsch in seiner ``Fundamentalaufgabe der Invariantentheorie'' zur Ableitung der Elementarkovarianten bei den Reihenentwicklungen. Mit Hilfe dieser Identitäten, die den Übergang von den nicht expliziten Determinantenausdrücken zum Matrizenprodukt vermitteln, läßt sich nun in der Tat nachweisen, daß mit der Matrizendarstellung die gewünschte explizite symbolische Darstellung gewonnen ist.

Für den Faltungsprozeß aber gilt ein ganz analoger Satz wie im ternären Gebiet, auf den sich auch analog die Reduktionssätze aufbauen lassen. Bezeichnet man in (1) die Zahl $\lambda$ als Defekt der Faltung, und Faltungen, für die $\sigma=\tau$, als kogredient, \emph{so lassen sich alle Faltungen erzeugen durch sukzessive Anwendung der kogredienten Faltungen vom Defekt 1}. Der Beweis dieses Satzes beruht wesentlich darauf, daß man die allgemeinsten Formen durch ``Normalformen'' ersetzen kann, d. h. durch Formen, die bei Anwendung aller invarianten Differentialoperationen identisch verschwinden. Im quaternären Gebiet hat diese letztere Tatsache Herr Mertens\srcfn{1)}{Über invariante Gebilde quaternärer Formen. Wiener Berichte. Bd. 98. 1889.} bewiesen; unsere Kenntnis der allgemeinsten Differentialoperationen -- die bekanntlich aus den Faltungsprozessen durch Ersetzen der Symbolreihen durch Differentialsymbole hervorgehen -- gestattet uns, unter Anwendung der Zerlegungsidentität den Mertensschen Beweis fast wörtlich auf das $n$-äre Gebiet zu übertragen.

